\section{Separable violations for every prescribed measurement pair}
\label{sec:allmub}

This section builds the separable violations in three steps. We first
compress each prescribed qudit measurement pair to a pair of binary
measurements on a plane (Section~\ref{sec:compression}). We then compute
the curvature of the gap for an explicit separable family on a qubit pair
(Section~\ref{sec:curvature}), and turn it into one finite witness that
works for every dimension and every prescribed pair
(Section~\ref{sec:witness}). The last subsection treats the case where both
local dimensions are at least three.

All logarithms in this section are natural.  For prescribed local
measurement pairs $(Q_A,R_A)$ and $(Q_B,R_B)$, write
\begin{equation}
 G(\rho)=I(A:B)_\rho-I(Q_A:Q_B)_\rho-I(R_A:R_B)_\rho.
 \label{eq:cqc-gap}
\end{equation}
Here $I(A:B)_\rho=S(\rho_A)+S(\rho_B)-S(\rho)$ is quantum mutual
information and the other terms are the Shannon mutual informations of
the measurement outcomes.  The CQC inequality asserts $G(\rho)\geq0$.
We take each measurement to be a rank-one projective measurement in an
orthonormal basis.  Two bases are mutually unbiased if all their squared
overlaps equal the reciprocal of the local dimension.

Our first result concerns a quantifier stronger than the existence of
one violating measurement pair.  Qian's all-MUB theorem establishes
separable violations for every prescribed Bob MUB pair when Bob's
dimension is at least $164$ \cite[Theorem~4.1]{qian}.
Counterexamples for \emph{some} MUB pair in the smaller dimensions are
also prior work \cite{qian}.  The following theorem applies to every
prescribed MUB pair throughout the qubit--qudit range $n\geq4$.

\begin{theorem}[Every prescribed MUB pair in qubit--qudit dimensions]
\label{thm:allmub-qubit}
 Let $n\geq4$.  For every prescribed ordered pair of mutually unbiased
 bases on $\mathbb C^2$, and every prescribed ordered pair of mutually
 unbiased bases on $\mathbb C^n$, there is a separable state $\rho$ of
 rank four such that
 \begin{equation}
  G(\rho)<-\frac{9}{20\,000\,000}.
  \label{eq:uniform-cqc-margin}
 \end{equation}
 For the same prescribed measurements, full-rank separable states with
 $G<0$ also exist.  The assertions hold with the two parties exchanged.
\end{theorem}

The construction uses a two-dimensional support on Bob, but the
measurements on that support are generally not mutually unbiased.
The state is evaluated with the original complete measurement tables;
no postselection is performed.

\subsection{Compression to two binary projective measurements}
\label{sec:compression}

Any pair of mutually unbiased qubit bases can be represented, after a
local unitary and outcome relabellings, by the Pauli observables $Z$ and
$X$.  We use this normal form on Alice.

\begin{lemma}[Binary compression of any MUB pair]
\label{lem:binary-compression}
 Let $Q=(q_j)_{j=0}^{n-1}$ and $R=(r_k)_{k=0}^{n-1}$ be mutually
 unbiased bases of $\mathbb C^n$, with $n\geq2$.  Choose $q_0\in Q$ and
 $r_0\in R$, and set
 \begin{equation}
  W=\operatorname{span}\{q_0,r_0\},\qquad P=P_W,
  \qquad B_1=2|q_0\rangle\langle q_0|-P,
  \qquad B_2=P-2|r_0\rangle\langle r_0|.
  \label{eq:mub-compression}
 \end{equation}
 On $W$, both $B_1$ and $B_2$ are traceless Hermitian involutions, and
 \begin{equation}
  \frac12\operatorname{Tr}_W(B_1B_2)=1-\frac2n.
  \label{eq:compressed-overlap}
 \end{equation}
 For every joint state supported on $\mathbb C^2\otimes W$,
 \begin{equation}
  I(Z_A:Q_B)=I(Z_A:B_1),\qquad
  I(X_A:R_B)=I(X_A:B_2).
  \label{eq:compression-mi}
 \end{equation}
\end{lemma}

\begin{proof}
 The vectors $q_0,r_0$ are linearly independent because
 $|\langle q_0,r_0\rangle|^2=1/n<1$.  The claimed trace properties
 follow directly from \eqref{eq:mub-compression}, and
 \[
  \frac12\operatorname{Tr}_W
  \bigl[(2|q_0\rangle\langle q_0|-P)
  (P-2|r_0\rangle\langle r_0|)\bigr]
  =1-2|\langle q_0,r_0\rangle|^2.
 \]
 For $j\ne0$, the projected vector $Pq_j$ lies in the single line
 $W\cap q_0^\perp$.  Thus the restricted $Q$ measurement is a
 projective measurement of $B_1$ followed by a fixed random subdivision
 of one binary outcome.  The binary outcome is, conversely, a
 deterministic function of the fine outcome.  The subdivision is
 independent of Alice conditional on that binary outcome, so the first
 equality in \eqref{eq:compression-mi} follows from the chain rule for
 mutual information.  The same argument applies to $R$; the sign in
 $B_2$ merely relabels the binary outcomes.
\end{proof}

The geometric compression in Lemma~\ref{lem:binary-compression} is
not new.  It is the one appearing in Qian's doubly degenerate stratum
\cite[Proposition~6.9]{qian}.  It is also O'Connor's ``logical qubit''
\cite{oconnor}: after rephasing so that $r_0=n^{-1/2}\sum_jq_j$, the
plane spanned by $q_0$ and the normalised sum of the remaining $Q$
vectors is exactly $W=\operatorname{span}\{q_0,r_0\}$.  We next use it
with a mixed-state perturbation different from the pure-ray mixture of
the earlier all-MUB tail theorem \cite[Theorem~4.1]{qian} and from the
entangled logical state of \cite{oconnor}.

\subsection{An explicit separable family and its exact curvature}
\label{sec:curvature}

We now forget the qudit and work with two binary observables on a qubit.
We write down a separable family and compute the second derivative of
its gap exactly.

In this subsection Bob is a qubit, with identity $\mathbf1_B$.  Let
$B_1,B_2$ be traceless Hermitian involutions.  Relabelling the outcomes
of one observable if necessary, assume
\begin{equation}
 c=\tfrac12\operatorname{Tr}(B_1B_2)\in[0,1].
 \label{eq:binary-overlap}
\end{equation}
On Alice put
\begin{equation}
 U=\frac{X+Z}{\sqrt2},\qquad
 V=\frac{Z-X}{\sqrt2},\qquad
 \alpha_r=\frac{\mathbf1_A+rU}{2},\qquad 0\leq r<1,
 \label{eq:alice-anchor}
\end{equation}
and define
\begin{equation}
 \sigma_r=\alpha_r\otimes\frac{\mathbf1_B}{2},\qquad
 \delta=\frac{Z\otimes B_1-X\otimes B_2}{4},\qquad
 \rho_{r,t}=\sigma_r+t\delta.
 \label{eq:separable-family}
\end{equation}
For this family the gap $G(r,c,t)$ means
$I(A:B)-I(Z:B_1)-I(X:B_2)$.

\begin{proposition}[Exact curvature and its threshold]
\label{prop:curvature-family}
 If $r+2|t|\leq1$, the operator $\rho_{r,t}$ is a separable density
 operator with marginals $\alpha_r$ and $\mathbf1_B/2$.  The gap is an
 even function of $t$, and
 \begin{equation}
 \begin{split}
  Q(r,c)
  &:=\left.\frac{\partial^2G(r,c,t)}{\partial t^2}\right|_{t=0}\\
  &=\frac{1-c}{1-r^2}
    +(1+c)\frac{\operatorname{arctanh}r}{r}
    -\frac{2}{1-r^2/2},
 \end{split}
 \label{eq:curvature-formula}
 \end{equation}
 with the continuous interpretation at $r=0$.
 There exists $r\in(0,1)$ with $Q(r,c)<0$ if and only if $c>1/2$.
 In particular, every $c>1/2$ yields separable violations for
 sufficiently small positive $r$ and then sufficiently small nonzero
 $t$.
\end{proposition}

\begin{proof}
 Both partial traces of $\delta$ vanish.  For an explicit separable
 decomposition let $\Pi_{j,b}=(\mathbf1_B+bB_j)/2$, $b\in\{-1,1\}$.
 Then
 \begin{equation}
  \rho_{r,t}=\frac14\sum_{b=\pm1}
  \left[(\alpha_r+btZ)\otimes\Pi_{1,b}
       +(\alpha_r-btX)\otimes\Pi_{2,b}\right].
  \label{eq:four-product-decomposition}
 \end{equation}
 Every Alice factor in this formula has trace one and a Bloch vector
 of norm at most $r+2|t|$.  The four summands are product density
 operators with possibly mixed Alice factors and pure Bob factors.
 This proves positivity, normalization and separability in the stated range.

 Bob has a unitary that sends both $B_1$ and $B_2$ to their negatives:
 take a Pauli observable perpendicular to their Bloch vectors.  Its
 conjugation sends $\rho_{r,t}$ to $\rho_{r,-t}$ and only relabels the
 binary measurement outcomes.  Hence the quantum mutual information
 and both classical mutual informations are even in $t$.

 The two classical tables agree up to an outcome relabelling.  Their
 entries are
 \begin{equation}
  p_{ab}(t)=\frac{1+a r/\sqrt2+abt}{4},
  \qquad a,b\in\{-1,1\}.
  \label{eq:binary-tables}
 \end{equation}
 Each has fixed marginals and mutual-information second derivative
 $1/(1-r^2/2)$ at zero.

 To compute the quantum term, let
 \[
  \Omega_\sigma(H)=\int_0^1\sigma^s H\sigma^{1-s}\,ds.
 \]
 The Hessian of $D(\sigma_r+t\delta\Vert\sigma_r)$ at zero is
 $\operatorname{Tr}[\delta\Omega_{\sigma_r}^{-1}(\delta)]$.
 This relative entropy is the quantum mutual information because the
 marginals are fixed.  In the eigenbasis of $U$,
 \begin{equation}
 \begin{split}
  \operatorname{Tr}[U\Omega_{\alpha_r}^{-1}(U)]
    &=\frac4{1-r^2},\\
  \operatorname{Tr}[V\Omega_{\alpha_r}^{-1}(V)]
    &=4\frac{\operatorname{arctanh}r}{r}.
 \end{split}
 \label{eq:alice-logmean}
 \end{equation}
 The first identity follows by adding the reciprocals of the two
 eigenvalues $(1\pm r)/2$; the second uses their logarithmic mean.
 The $U$ and $V$ parts are orthogonal for this form.  Since
 \[
  \delta=\frac{U\otimes(B_1-B_2)+V\otimes(B_1+B_2)}{4\sqrt2},
  \qquad
  \operatorname{Tr}\bigl[(B_1\mp B_2)^2\bigr]=4(1\mp c),
 \]
 and $\Omega_{\sigma_r}^{-1}$ contributes a factor two on the
 maximally mixed Bob factor, the quantum second derivative is the
 sum of the first two terms in \eqref{eq:curvature-formula}.
 Subtracting the two classical derivatives proves that formula.

 Finally, its absolutely convergent power series for $|r|<1$ is
 \begin{equation}
  Q(r,c)=\sum_{k\geq1}
   \left[(1-c)+\frac{1+c}{2k+1}-2^{1-k}\right]r^{2k}.
  \label{eq:curvature-series}
 \end{equation}
 The coefficient at $r^2$ is $(1-2c)/3$, so $c>1/2$ gives a negative
 value for all sufficiently small positive $r$.  At $c=1/2$ that
 coefficient vanishes, while every coefficient for $k\geq2$ is
 strictly positive:
 \[
  \frac12+\frac{3}{2(2k+1)}-2^{1-k}>0.
 \]
 Every coefficient decreases as $c$ increases, so all coefficients
 are nonnegative for $c\leq1/2$.  This proves the threshold assertion.
 The gap vanishes at $t=0$ and is even, so a strictly negative second
 derivative gives a strictly negative gap for sufficiently small
 nonzero $t$, within the separable range already proved.
\end{proof}

The threshold in Proposition~\ref{prop:curvature-family} is exact for
the product-state Hessian of the displayed family.  It does not assert
a global CQC theorem for all two-qubit states when $c\leq1/2$.

For direct evaluation, the four eigenvalues of
\eqref{eq:separable-family} are
\begin{equation}
 \lambda_{\eta,\xi}(r,c,t)
 =\frac{1+\eta t\sqrt{1-c}
             +\xi\sqrt{r^2+(1+c)t^2}}4,
 \qquad\eta,\xi\in\{-1,1\}.
 \label{eq:four-eigenvalues}
\end{equation}
Indeed, for $c<1$ set
$C=(B_1-B_2)/\sqrt{2(1-c)}$ and
$D=(B_1+B_2)/\sqrt{2(1+c)}$.  These are orthogonal Pauli
observables, and
\[
 4\rho_{r,t}=\mathbf1+rU\otimes\mathbf1_B
   +t\sqrt{1-c}\,U\otimes C+t\sqrt{1+c}\,V\otimes D.
\]
The $U\otimes C$ term commutes with the other two nonconstant terms,
and those two anticommute.  This proves
\eqref{eq:four-eigenvalues}, including $c=1$ by continuity.

\subsection{A uniform finite witness}
\label{sec:witness}

The curvature formula shows that violations exist for small parameters.
We now bound the higher-order terms by rational estimates, valid for any
admissible parameters, and then fix one parameter choice that works for
every overlap $c\geq3/5$.

\begin{lemma}[Quartic remainder bound]
\label{lem:remainder}
 Let $0<r<1$, $t>0$ and $r+2t<1$, and let $c\in[0,1]$.  Let
 $\bar N$ and $\bar A(r)$ be numbers with
 \begin{equation}
  \bar N\geq\|\delta\|=N_c,\quad N_c^2=\frac{1+\sqrt{1-c^2}}8,
  \qquad
  \bar A(r)\geq\frac{\operatorname{arctanh}r}{r},
  \label{eq:delta-norms}
 \end{equation}
 and put
 \[
  L=\frac{1-r}4-t\bar N,\qquad
  \bar Q=\frac{1-c}{1-r^2}+(1+c)\bar A(r)-\frac{2}{1-r^2/2}.
 \]
 If $L>0$, then $\rho_{r,t}$ is separable and positive definite on its
 four-dimensional support, and
 \begin{equation}
  G(r,c,t)\leq\frac{\bar Q}{2}t^2+\frac{N_c^2}{24L^3}t^4.
  \label{eq:finite-gap-bound}
 \end{equation}
\end{lemma}

\begin{proof}
 Separability is \eqref{eq:four-product-decomposition}, since
 $r+2t<1$.  Direct Pauli algebra gives $\|\delta\|_F^2=1/2$ and the
 operator norm in \eqref{eq:delta-norms}.  On the support, the smallest
 eigenvalue of $\sigma_r$ is $(1-r)/4$, so by Weyl's inequality
 $\lambda_{\min}(\rho_{r,s})\geq L>0$ for $|s|\leq t$.
 Let $f(s)=\operatorname{Tr}(\rho_{r,s}\log\rho_{r,s})$.
 Differentiating the resolvent representation for the second
 derivative gives
 \begin{equation}
  f^{(4)}(s)=6\int_0^\infty
   \operatorname{Tr}\left[
    \bigl(\delta(\rho_{r,s}+u\mathbf1)^{-1}\bigr)^4
   \right]du.
  \label{eq:fourth-resolvent}
 \end{equation}
 If $R=(\rho_{r,s}+u\mathbf1)^{-1}$, the integrand is
 $\operatorname{Tr}\bigl[(R^{1/2}\delta R^{1/2})^4\bigr]$ and is nonnegative.
 Applying $\operatorname{Tr}A^4\leq\|A\|^2\|A\|_F^2$ to this
 Hermitian matrix yields
 \begin{equation}
  0\leq f^{(4)}(s)\leq
   \frac{2\|\delta\|^2\|\delta\|_F^2}{\lambda_{\min}(\rho_{r,s})^3}
   \leq\frac{N_c^2}{L^3}.
 \label{eq:fourth-bound}
 \end{equation}
 The quantum mutual information equals $f(s)-f(0)$, and is even in
 $s$.  Each classical mutual information $J(s)$ is also even.  Because
 the marginals of \eqref{eq:binary-tables} are fixed and its cells are
 affine in $s$,
 \[
  J^{(4)}(s)=2\sum_{a,b}\frac{(p'_{ab}(s))^4}{p_{ab}(s)^3}
  \geq0,
 \]
 where all cells $(1\pm r/\sqrt2\pm s)/4$ are positive because
 $r/\sqrt2+t<1$, which follows from $r+2t<1$.  Thus $J(t)\geq J''(0)t^2/2$.  Since the second derivative
 of the gap at zero is $Q(r,c)\leq\bar Q$, Taylor's theorem with the
 Lagrange remainder gives \eqref{eq:finite-gap-bound}.
\end{proof}

\begin{lemma}[Uniform strictly negative gap]
\label{lem:finite-witness}
 For every $c\in[3/5,1]$, the state
 $\rho_{61/200,\,421/20000}$ in \eqref{eq:separable-family} is
 separable, is positive definite on its four-dimensional support, and
 satisfies \eqref{eq:uniform-cqc-margin}.
\end{lemma}

\begin{proof}
 Apply Lemma~\ref{lem:remainder} with $(r,t)=(61/200,421/20000)$,
 \begin{equation}
  \bar A(r)=\sum_{k=0}^{40}\frac{r^{2k}}{2k+1}
            +\frac{r^{82}}{83(1-r^2)},
  \label{eq:artanh-upper}
 \end{equation}
 which bounds $\operatorname{arctanh}(r)/r$ because every omitted
 coefficient is at most $1/83$, and with $\bar N$ a rational number at
 least $\sqrt{N_c^2}$, rounded upwards at denominator $10^{12}$ (the
 same rounding is applied to $\sqrt{1-c^2}$ inside $N_c^2$).  The right
 side of \eqref{eq:finite-gap-bound} decreases in $c$: its $c$-derivative
 through $\bar Q$ is $\bar A(r)-1/(1-r^2)<0$, and $N_c^2$ decreases
 while $L$ increases.  It therefore suffices to evaluate it at $c=3/5$,
 where $N_c^2=9/40$; the script uses a rational upper bound for
 $N_c^2$ (rounded upwards as above), which can only weaken the bound.  This is a finite computation in exact rational
 arithmetic.  It gives a value below $-9/20\,000\,000$; the script
 \texttt{scout\_cqc/certified\_margin.py} performs it, together with
 the monotonicity inequality and a check on a grid of $41$ values of
 $c$, with no floating-point arithmetic.
\end{proof}

\begin{remark}[A hand-checkable cruder constant]
\label{rem:old-constant}
 The same lemma at $(r,t)=(1/3,1/100)$ can be checked by hand and gives,
 for every $c\in[3/5,1]$, the weaker margin
 $-17117/114\,750\,000\,000\approx-1.49\times10^{-7}$
 of an earlier version of this work (the bounds below use $c\geq3/5$).  At $r=1/3$,
 \eqref{eq:curvature-formula} gives
 $Q(1/3,c)=\tfrac98(1-c)+\tfrac32(1+c)\log2-\tfrac{36}{17}$, which
 decreases with $c$.  The series bound
 \begin{equation}
  \log2=2\sum_{k\geq0}\frac1{(2k+1)3^{2k+1}}
   <\frac23+\frac2{81}+\frac1{540}=\frac{1123}{1620}
 \label{eq:log-two-upper}
 \end{equation}
 follows by bounding $1/(2k+1)$ by $1/5$ for $k\geq2$; it is strict for
 $k>2$.  Hence
 \begin{equation}
  Q(1/3,c)\leq\frac{12}{5}\log2-\frac{567}{340}<-\frac{181}{45900},
  \label{eq:uniform-negative-curvature}
 \end{equation}
 where $(12/5)(1123/1620)-567/340=-181/45900$ exactly, so strictness
 comes solely from \eqref{eq:log-two-upper}.  Using the cruder bounds
 $\|\delta\|\leq1/2$ and $\|\delta\|^2\leq9/40$ for $|s|\leq1/12$ gives
 $\lambda_{\min}\geq1/6-|s|/2\geq1/8$ and $f^{(4)}\leq576/5$, so
 $G(1/3,c,t)<-181t^2/91800+(24/5)t^4$, which equals
 $-17117/114\,750\,000\,000$ at $t=1/100$.  The rational evaluation of
 Lemma~\ref{lem:remainder} at this point gives about
 $-1.83\times10^{-7}$, so the earlier constant also follows from it.
\end{remark}

\paragraph{Margins for each dimension.}
Optimising $(r,t)$ separately for each overlap gives larger certified
margins.  Table~\ref{tab:margins} lists them; each entry is proved by the
same exact rational evaluation of Lemma~\ref{lem:remainder} in
\texttt{certified\_margin.py}.  For $n=4$ the overlap is $c=1$
(see the proof below).  Because the bound decreases in $c$, the entry
for $n=12$ also holds for every $n\geq13$, and the $c=1$ entry is the
limit as $n\to\infty$.

\begin{table}[ht]
\centering\small
\caption{Certified separable CQC violations for each Bob dimension $n$,
with $c=1-2/n$ for $n\geq5$ and $c=1$ for $n=4$.  Each margin is a
proved strict upper bound on the gap $G$, in nats.}
\label{tab:margins}
\begin{tabular}{rccc}
\toprule
$n$ & $c$ & $(r,t)$ & certified $G<$\\
\midrule
4 & $1$ & $(113/200,\ 2851/50000)$ & $-1.058\times10^{-4}$\\
5 & $3/5$ & $(61/200,\ 421/20000)$ & $-4.572\times10^{-7}$\\
6 & $2/3$ & $(19/50,\ 2881/100000)$ & $-2.522\times10^{-6}$\\
7 & $5/7$ & $(421/1000,\ 3301/100000)$ & $-5.673\times10^{-6}$\\
8 & $3/4$ & $(223/500,\ 3569/100000)$ & $-9.249\times10^{-6}$\\
9 & $7/9$ & $(93/200,\ 1879/50000)$ & $-1.288\times10^{-5}$\\
10 & $4/5$ & $(239/500,\ 39/1000)$ & $-1.639\times10^{-5}$\\
11 & $9/11$ & $(61/125,\ 1003/25000)$ & $-1.971\times10^{-5}$\\
12 & $5/6$ & $(62/125,\ 4103/100000)$ & $-2.281\times10^{-5}$\\
\bottomrule
\end{tabular}
\end{table}

\paragraph{How large are the true violations? (Numerical, not a proof.)}
The quartic remainder bound is loose at larger $t$, so the true gaps are
larger than the certified ones.  Direct numerical evaluation of the gap in
the same family, optimised over $(r,t)$, gives the values in
Table~\ref{tab:numeric}.  At $n=5$ the best value in the family is about
twelve times the value at the hand-checkable point $(1/3,1/100)$.  These
numbers are illustrations only; they are not used in any proof.

\begin{table}[ht]
\centering\small
\caption{Numerical (non-proof) gaps in the family
\eqref{eq:separable-family}, in nats.}
\label{tab:numeric}
\begin{tabular}{rcccc}
\toprule
$n$ & $c$ & best $G$ in the family & at $(r,t)\approx$ & $G$ at $(1/3,1/100)$\\
\midrule
4 & $1$ & $-2.43\times10^{-3}$ & $(0.573,\ 0.214)$ & $-1.91\times10^{-6}$\\
5 & $0.6$ & $-2.35\times10^{-6}$ & $(0.329,\ 0.048)$ & $-2.00\times10^{-7}$\\
6 & $0.667$ & $-2.03\times10^{-5}$ & $(0.421,\ 0.081)$ & $-4.85\times10^{-7}$\\
7 & $0.714$ & $-6.15\times10^{-5}$ & $(0.475,\ 0.107)$ & $-6.88\times10^{-7}$\\
8 & $0.75$ & $-1.25\times10^{-4}$ & $(0.511,\ 0.127)$ & $-8.41\times10^{-7}$\\
\bottomrule
\end{tabular}
\end{table}


We prove Theorem~\ref{thm:allmub-qubit} first for $n\geq5$, where the
compression of Lemma~\ref{lem:binary-compression} has $c\geq3/5$, and
then for $n=4$, where it has only $c=1/2$ and a different plane is
needed.

\begin{proof}[Proof of Theorem~\ref{thm:allmub-qubit}]
 For $n\geq5$, use Lemma~\ref{lem:binary-compression}; its overlap
 is $c=1-2/n\geq3/5$.  Embed the state of
 Lemma~\ref{lem:finite-witness} in $\mathbb C^2\otimes W$.
 Written on the original space it is
 \begin{equation}
  \rho=\frac14\left[
   \left(\mathbf1_A+\frac{61(X+Z)}{200\sqrt2}\right)\otimes P
   +\frac{421}{20000}(Z\otimes B_1-X\otimes B_2)
  \right].
  \label{eq:explicit-all-mub-state}
 \end{equation}
 It has rank four.  Its separability follows from
 \eqref{eq:four-product-decomposition}, with
 $\Pi_{j,b}=(P+bB_j)/2$, and the mutual-information equalities
 \eqref{eq:compression-mi} transfer the strict gap to the original
 measurement tables.

 For $n=4$, write $Q$ and $R$ also for the unitary matrices whose
 columns are the prescribed basis vectors.  The matrix $2Q^*R$ is a
 complex Hadamard matrix of order four.  By the exhaustive
 classification at this order \cite[Section~5.4, equations~(62)--(65)]{tadej},
 $Q^*R=D_1P_1H(a)P_2D_2$ for diagonal unitaries $D_i$, permutation
 matrices $P_i$, and
 \begin{equation}
 H(a)=\frac12\begin{pmatrix}
  1&1&1&1\\
  1&a&-1&-a\\
  1&-1&1&-1\\
  1&-a&-1&a
 \end{pmatrix},\qquad |a|=1.
 \label{eq:hadamard-four}
 \end{equation}
 Put $V=QD_1P_1$.  Then $Q=V(D_1P_1)^{-1}$ and $R=VH(a)P_2D_2$, so the
 rays of $Q$ are the columns of $V$ and the rays of $R$ are the columns
 of $VH(a)$, up to order and phase.  Set
 \[
  u=V\frac{(1,0,1,0)^{\mathsf T}}{\sqrt2},\qquad
  v=V\frac{(0,1,0,-1)^{\mathsf T}}{\sqrt2}.
 \]
 The left factor $D_1$ must be kept in $V$: with $V$ replaced by
 $QP_1$ the vectors are generally different and the construction
 fails.  In the frame $V$, the two vectors have disjoint supports in the
 computational basis, and also in the basis of columns of $H(a)$,
 because
 \[
  H(a)^*V^*u=\frac{(1,0,1,0)^{\mathsf T}}{\sqrt2},\qquad
  H(a)^*V^*v=\frac{(0,\bar a,0,-\bar a)^{\mathsf T}}{\sqrt2}.
 \]
 On $W=\operatorname{span}\{u,v\}$, both fine measurements $Q$ and $R$
 are consequently the same binary projective measurement, followed by
 independent subdivisions of its outcomes.  Take
 $B_1=B_2=|u\rangle\langle u|-|v\rangle\langle v|$;
 now $c=1$, and Lemma~\ref{lem:finite-witness} again gives
 \eqref{eq:uniform-cqc-margin}.  The state
 \eqref{eq:explicit-all-mub-state}, with $P=|u\rangle\langle u|+|v\rangle\langle v|$
 and these $B_1,B_2$, is already written for the originally prescribed
 pair.  The accompanying numerical checks include random instances of
 $Q$, $D_1,P_1,D_2,P_2$ and $a$.

 Finally, for either construction set
 \[
  \rho_\varepsilon=(1-\varepsilon)\rho
                  +\varepsilon\frac{\mathbf1_{2n}}{2n}.
 \]
 For every $0<\varepsilon<1$ this is full rank and separable.
 Continuity of the finite-dimensional quantum and classical entropies
 preserves $G(\rho_\varepsilon)<0$ for sufficiently small
 $\varepsilon>0$.  This last assertion does not require, or assert, a
 noise tolerance uniform in $n$.  Exchanging the parties proves the
 transposed statement.
\end{proof}

At $n=4$ the same disjoint-support vectors also give the familiar
rank-two double-readout witness
$\tfrac12|0,u\rangle\langle0,u|+
\tfrac12|+,v\rangle\langle+,v|$ from the existing counterexample
mechanism \cite[Appendix~A.3]{qian}.  The proof above instead uses
the same finite parameters and rational margin in every dimension.

\subsection{When both local dimensions are at least three}

The locally orthogonal shared-label mechanism of the original
counterexamples \cite{wangcounter,qian} also extends to arbitrary
prescribed MUB pairs.  For equal dimensions $d\geq3$, an adapted
every-prescribed-pair construction was stated earlier by O'Connor
\cite{oconnor}, using the compression of
Lemma~\ref{lem:binary-compression} and an entangled logical state; its
author flagged it as not independently checked.  The theorem below
gives a separable version, allows unequal dimensions, and adds a
margin independent of the dimensions, with a complete proof.  We
record it separately from the qubit curvature construction.

\begin{theorem}[Arbitrary prescribed bases in higher local dimensions]
\label{thm:allmub-higher}
 Let $m,n\geq3$.  For every prescribed MUB pair on $\mathbb C^m$
 and every prescribed MUB pair on $\mathbb C^n$, there is a
 rank-two separable state $\tau$, diagonal in a local product basis,
 such that
 \begin{equation}
  I(Q_A:Q_B)_\tau+I(R_A:R_B)_\tau-I(A:B)_\tau
  \geq\frac1{2\pi^4}.
  \label{eq:higher-uniform-gap}
 \end{equation}
 Full-rank separable states with strictly positive excess exist for
 the same prescribed measurements.
\end{theorem}

\begin{proof}
 On a local space of dimension $d\geq3$, select three distinct vectors
 $q_0,q_1,q_2$ of $Q$ and define
 \[
  u=q_0,\qquad
  v_\theta=\frac{q_1+e^{i\theta}q_2}{\sqrt2}.
 \]
 The $R$ law of $u$ is uniform, $p(k)=1/d$.  Mutual unbiasedness
 gives phases $z_k$, with $|z_k|=1$, for which the $R$ law of
 $v_\theta$ is
 \[
  q_\theta(k)=\frac{1+\operatorname{Re}(e^{i\theta}z_k)}d.
 \]
 Averaging over the phase shows
 \begin{equation}
  \frac1{2\pi}\int_0^{2\pi}\|q_\theta-p\|_1\,d\theta
  =\frac2\pi.
  \label{eq:phase-average}
 \end{equation}
 Choose a phase attaining at least this average, independently on
 Alice and Bob.  Set
 \[
  \tau=\tfrac12|u_Au_B\rangle\langle u_Au_B|
       +\tfrac12|v_Av_B\rangle\langle v_Av_B|.
 \]
 The product rays are locally orthogonal, so $\tau$ and both of its
 marginals have nonzero spectrum $(1/2,1/2)$.  It follows that
 $I(A:B)_\tau=\log2$.  Each $Q$ measurement identifies the latent
 mixture label exactly, and the local outcomes are independent
 conditional on that label.  Hence $I(Q_A:Q_B)_\tau=\log2$ as well.

 For the $R$ outcomes, write $p_A,p_B$ for the conditional uniform
 laws and $q_A,q_B$ for the other conditional laws.  Their joint law
 is $P=\tfrac12p_A\otimes p_B+\tfrac12q_A\otimes q_B$, and
 \[
  P-P_A\otimes P_B
  =\tfrac14(p_A-q_A)\otimes(p_B-q_B).
 \]
 Therefore \eqref{eq:phase-average} gives
 \[
  \|P-P_A\otimes P_B\|_1
   =\tfrac14\|p_A-q_A\|_1\|p_B-q_B\|_1
   \geq\frac1{\pi^2}.
 \]
 Pinsker's inequality in natural units now yields
 \[
  I(R_A:R_B)_\tau
    =D(P\Vert P_A\otimes P_B)
    \geq\tfrac12\|P-P_A\otimes P_B\|_1^2
    \geq\frac1{2\pi^4},
 \]
 which is precisely the excess in \eqref{eq:higher-uniform-gap}.
 The full-rank assertion follows by mixing with sufficiently little
 maximally mixed noise and using continuity, as above.
\end{proof}

For a finite construction without optimizing a continuous phase,
one may choose only between $\theta=0$ and $\theta=\pi/2$ on each
party.  The elementary inequality $|\operatorname{Re}z|+
|\operatorname{Im}z|\geq1$ for $|z|=1$ ensures that one choice gives
$\|q_\theta-p\|_1\geq1/2$.  The same proof then gives the uniform
excess $1/512$.  The stronger bound \eqref{eq:higher-uniform-gap}
uses the phase average.

For completeness, the required form of Pinsker's inequality can be
proved without additional matrix inequalities.  Coarsen two
probability laws $P,Q$ into the event $\{P>Q\}$ and its complement.
The log-sum inequality bounds $D(P\Vert Q)$ below by the binary
relative entropy $d(a\Vert b)$, with
$a-b=\|P-Q\|_1/2$.  As a function of $a$, this binary relative
entropy vanishes together with its first derivative at $a=b$, and
has second derivative $1/[a(1-a)]\geq4$.  Thus
$d(a\Vert b)\geq2(a-b)^2=\|P-Q\|_1^2/2$.

Theorems~\ref{thm:allmub-qubit} and~\ref{thm:allmub-higher}
make the failure independent of which MUB pairs are prescribed in
the covered dimensions.  They do not decide the remaining global
$2\times3$ CQC problem: the two-vector compression there has
$c=1/3$, below the threshold of
Proposition~\ref{prop:curvature-family}.
